\clearpage
\renewcommand{\sectioncolor}{proof}
\section{How to use these three things together}
\markboth{How to use this}{}

\begin{center}
\begin{tikzpicture}[font=\footnotesize,
 b/.style={rounded corners=4pt,draw=#1,line width=1.2pt,fill=#1!8,inner sep=9pt,
           text width=4.9cm,align=left,minimum height=3.6cm},
 ar/.style={-{Stealth[length=5.5pt]},line width=1.2pt,draw=slate!70}]
 \node[b=bookcol] (a) at (-5.6,0) {\textbf{\color{bookcol}THE BOOKS}\\[4pt]
  \bk{Maths\_Proof\_Skill/}\\[4pt]{\scriptsize 8 books, 2{,}164 pp.\\
  You read \textbf{134} of them: Devlin §3 (25 pp) and Thurston (17 pp), then dip into
  Diedrichs \& Lovett for drills.}};
 \node[b=proof] (b) at (0,0) {\textbf{\color{proof}THIS DOCUMENT}\\[4pt]
  {\scriptsize Explains the notebook cell by cell, and pulls the \emph{method} out of the books ---
  definitions, the nine steps for finding a proof, the six techniques, and the trap.}\\[4pt]
  {\scriptsize Read it beside the notebook.}};
 \node[b=nano] (c) at (5.6,0) {\textbf{\color{nano}THE NOTEBOOK}\\[4pt]
  \bk{Proof\_Techniques\_01\_...ipynb}\\[4pt]{\scriptsize 22 cells. Where you \emph{do} it ---
  run the cells, hunt the counterexamples, and watch a CAS mislead you.}};
 \draw[ar] (a)--(b); \draw[ar] (b)--(c);
\end{tikzpicture}
\end{center}
\vspace{2pt}
\begin{center}\begin{minipage}{0.93\textwidth}\footnotesize\color{slate}
\textbf{Figure 1.}\; The books supply the method, this document supplies the map, the notebook is
where you practise. None of the three works well alone.
\end{minipage}\end{center}

\section{What a proof is --- three definitions, from three of your books}
\markboth{What a proof is}{}

\begin{bookbox}[title={Devlin, \textit{Introduction to Mathematical Thinking}, §3.1 (p.~117 book, \bk{01\_Start\_Here})}]
\begin{quote}\itshape
``Proofs are constructed for two main purposes: to establish truth and to communicate to others\ldots\
a proof of a statement must explain \textbf{why} that statement is true.''
\end{quote}
Note the second purpose. A proof that convinces only you is half a proof.
\end{bookbox}

\begin{bookbox}[title={Diedrichs \& Lovett, \textit{Transition to Advanced Mathematics}, Definition 2.1.15 (\bk{02\_Main\_Reference})}]
\begin{quote}\itshape
``A proof of a mathematical statement is a \textbf{sound argument} that establishes the truth of the
conclusion from the truth of the stated hypotheses.''
\end{quote}
Sharper and more formal. Notice what it does \emph{not} say: nothing about examples, computation or
evidence.
\end{bookbox}

\begin{bookbox}[title={Thurston, \textit{On Proof and Progress in Mathematics}, p.~2 (\bk{01\_Start\_Here})}]
\begin{quote}\itshape
``How do mathematicians advance human \textbf{understanding} of mathematics? \ldots\ what we are
doing is finding ways for people to understand and think about mathematics.''
\end{quote}
Thurston reframes the whole activity. He then gives the four-colour theorem as his example: it was
proved by massive computation, and the controversy
\begin{quote}\itshape
``had little to do with doubt people had as to the veracity of the theorem or the correctness of the
proof. Rather, it reflected a continuing desire for \textbf{human understanding} of a proof, in
addition to knowledge that the theorem is true.''
\end{quote}
\textbf{That is your situation exactly}, every time Sage prints \bk{True}.
\end{bookbox}

\begin{techbox}[title={Putting the three together}]
A proof is
\begin{enumerate}[label=\textbf{\arabic*.}]
\item a \textbf{sound argument} (Diedrichs \& Lovett) --- each step follows by a rule of logic from
 definitions, axioms, or theorems already proved;
\item that \textbf{explains why} (Devlin) --- not merely that;
\item and \textbf{communicates} that explanation to a reader (Devlin, Thurston).
\end{enumerate}
A computation can satisfy none of these and still be correct. That is the whole difficulty.
\end{techbox}
